\documentclass{beamer}
  \usepackage[utf8]{inputenc}
  \usepackage{hyperref}
  \usepackage{xspace}
  \usepackage{xcolor}
  \usepackage{graphicx}
  \usepackage{Sweave}
  \usepackage{multicol}
  
  \definecolor{links}{HTML}{2A1B81}
  \hypersetup{colorlinks,linkcolor=,urlcolor=links}  
  
  %Beamer properties
  \usetheme{Warsaw} %{Boadilla}
  \useoutertheme{miniframes}
  \usecolortheme{seahorse}
  \usecolortheme{rose}
  \setbeamerfont{footnote}{size=\tiny}

  % Add R Logo
  \definecolor{Rcolor}{RGB}{150,160,190}

  \newcommand{\Rx}{\fontsize{10pt}{12pt}\selectfont
  	\raisebox{.3em}{\hspace{1.2em}%
  		\llap{\resizebox{1.09em}{.5em}{\color{black}$\bigcirc$}}%
  		\llap{\resizebox{1.199em}{.55em}{\color{darkgray}$\bigcirc$}}%
  		\llap{\resizebox{1.19em}{.52em}{\color{gray!50}$\bigcirc$}}%
		\llap{\resizebox{1.1em}{.5em}{\color{gray}$\bigcirc$}}%
		\llap{\resizebox{1.25em}{.55em}{\color{gray}$\bigcirc$}}%
	}%
	\hspace{-.85em}%
	\textbf{%
		%\resizebox{.55em}{1.5ex}{\textcolor{black!80}{\textsf{R}}}%
		\textcolor{black}{\textsf{R}}%
		\hspace{-.025em}\raisebox{.01em}{\llap{\textcolor{Rcolor}{\textsf{R}}}}%
	}}%

  \newbox\rbox
  \savebox\rbox{\scalebox{0.1}{\Rx}}
  \makeatletter
  \def\R{\scalebox{\f@size}{\usebox\rbox}\xspace}
  \makeatletter
  
  %S4 representations
  \newenvironment<>{rsdmx}[1]{%
  \DefineVerbatimEnvironment{Soutput}{Verbatim}{fontsize=\tiny,formatcom=\color{black}}
  \setbeamerfont{block title}{size=\normalsize}
  \setbeamerfont{block body}{size=\scriptsize}
  \setbeamercolor{block title}{fg=white,bg=gray!90!black}%
  \begin{block}#2{#1 \R representation}}{\end{block}}
  
  %R code
  \DefineVerbatimEnvironment{Sinput}{Verbatim}{fontsize=\scriptsize,formatcom=\color{blue}}
  \DefineVerbatimEnvironment{Soutput}{Verbatim}{fontsize=\tiny,formatcom=\color{gray}}

  %RSDMX Presentation%
  %------------------%

  \title{rsdmx - Tools for reading SDMX data and metadata documents in \R}
  \author[E. Blondel]{Emmanuel Blondel \\ \texttt{emmanuel.blondel1@gmail.com}}
  \institute{CEO - International Consultant}
  \date{September 27, 2015}

\AtBeginSection[]
{
  \begin{frame}<beamer>
    \frametitle{Outline}
    \tableofcontents[currentsection]
  \end{frame}
}

%\setbeamertemplate{section in toc}{\protect\block{Research Question~\inserttocsectionnumber}\inserttocsection\protect\endblock}

\begin{document}

  \begin{frame}
  	\titlepage
  \end{frame}

\section{Introduction}
  
  %Introducing SDMX%
  %----------------%
  \begin{frame}
  	\frametitle{SDMX}
  	\framesubtitle{Statistical Data and Metadata Exchange (SDMX)}
  	Joint initiative created in 2001 by international \& regional institutions \footnotemark
  	\begin{itemize}
		\item Promote and develop standards and guidelines for the exchange and sharing of statistical \textbf{data} and \textbf{metadata}
			\begin{itemize}
				\item Definition of an abstract information model
				\item Development of standard formats
				\item Design of web-services architectures and tools 
			\end{itemize}
		\item Continuous process of improving the exchange of statistical data \& metadata
			\begin{itemize}
				\item an evolving set of specifications: SDMX 1.0, 2.0, 2.1
				\item a variety of formats: SDMX-ML, SDMX-EDI, SDMX-JSON
				\item a variety of service architectures: SOAP, REST
			\end{itemize}	
		\item Main SDMX format used across institutions: SDMX-ML
  		
  	\end{itemize}
  	
  	 \footnotetext[1]{Bank for International Settlements (BIS), European Central Bank (ECB), Statistical Office of the European Union (EUROSTAT), International Monetary Fund (IMF), Organization for Economic Co-operation and Development (OECD),United Nations (UN) and World Bank}
  \end{frame}

  %Conciliating SDMX and R%
  %-----------------------%
  \begin{frame}
  	\frametitle{Motivation}
  	\framesubtitle{Conciliating SDMX and R}
  	\begin{itemize}
		\item Need of Interoperability between statistical systems, formats and tools
		\item Need to co-analyse and process statistical data
			\begin{itemize}
				\item from a variety of domains (demography, socio-economics, health, environment, agriculture, fishery, etc.)
				\item from scattered data providers (national, regional \& international institutions)
				\item by a growing range of actors (\textit{e.g.} government institutions, statistical institutes, non-profit organizations, universities, research centers, companies)
			\end{itemize}						
		\item Fundamental and growing role of \R as plateform for statistical computing
		\item ...and the need of tools to facilitate reading of SDMX data and metadata in \R 
  	
  	\end{itemize}
  \end{frame}
  
  %rsdmx package introduction - 1%
  %------------------------------%
  \begin{frame}
  	\frametitle{rsdmx}
  	\framesubtitle{Introduction}

  	\begin{itemize}
		\item First initiative to read SDMX in R made available to the R community		
		\item Set of tools to read SDMX-ML data and metadata documents				
  	\end{itemize}
	  	
  \end{frame}

  %rsdmx package introduction - 2%
  %------------------------------%
  \begin{frame}
  	\frametitle{rsdmx}
  	\framesubtitle{Introduction}
  	
  	\begin{itemize}
		\item Generic SDMX abstraction library for R
		\item Read SDMX documents in a flexible way:
			\begin{itemize}
				\item Support for SDMX-ML 1.0, 2.0 and 2.1 standard formats
				\item Support for \textit{remote} or \textit{local} sources,
				\item No restriction to the SDMX web-services standard specifications (for \textit{remote} sources)
			\end{itemize}
		\item Variety of SDMX documents:
			\begin{itemize}
				\item \textbf{Data} (generic, compact, structure-specific, etc.)
				\item \textbf{Metadata} (Data structure definition - DSD, Codelists, Concepts, etc.)
			\end{itemize}						
  	\end{itemize}
	  	
  \end{frame}
  
  
  %rsdmx package introduction - 3%
  %------------------------------%
  \begin{frame}
  	\frametitle{rsdmx}
  	\framesubtitle{Introduction}
  	
  	\begin{multicols}{2}
  	  	
  	\begin{itemize}
		\item a single \texttt{readSDMX} function, with a \textbf{large bandwidth of use}:
			\begin{itemize}
				\item "raw" approach (read from \textit{url} or \textit{file})
				\item "helping" approach (read from a list of well-known service providers, with no need to specify the entire request)
			\end{itemize}
		\item a set of generic methods to convert SDMX data into tabular data (\texttt{data.frame})
  	\end{itemize}
  	
  	\columnbreak
	\hspace{2cm}
	\raggedleft
	\includegraphics[width=6cm]{./img/rsdmx.png}

	\end{multicols}
	  	
  \end{frame}
 
\section{Architecture of rsdmx} 
  
  %rsdmx package - S4 model%
  %------------------------%
  \begin{frame}[containsverbatim]
  	\frametitle{rsdmx}
  	\framesubtitle{Architecture - Object-oriented model}
  	\begin{itemize}
  		\item SDMX model represented in \R with \texttt{S4} classes and methods
  		\item In \texttt{S4} modelling, a \texttt{class} is made of \textit{slots} (properties)
  		\item the general structure of SDMX-ML document is represented with an \texttt{SDMX} abstract class
		\item each SDMX-ML document has a corresponding \R \texttt{SDMX*} class that extends the \texttt{SDMX} class	
  	\end{itemize}
  	
\begin{rsdmx}{SDMX}
\begin{multicols}{2}

\begin{Schunk}
\begin{Soutput}
Class "SDMX" [package "rsdmx"]
Slots:
                                                                  
Name:               xmlObj        schema        header
Class: XMLInternalDocument    SDMXSchema    SDMXHeader
				
Known Subclasses:
"SDMXGenericData", "SDMXCompactData", "SDMXMessageGroup",
"SDMXConcepts", "SDMXCodelists", "SDMXDataStructures",
"SDMXDataStructureDefinition"
\end{Soutput}
\end{Schunk} 

\columnbreak
\hspace{2cm}
\raggedleft
	\includegraphics[width=3.5cm]{img/SDMX.png}


\end{multicols}
\end{rsdmx}
  \end{frame}
  
  %rsdmx package - supported messages %
  %-----------------------------------%
  \begin{frame}
  	\frametitle{rsdmx}
  	\framesubtitle{Architecture - Supported SDMX-ML documents}
	\begin{itemize}
		\item \textit{Structure} types, \textit{i.e.} the elements that define the data structure, including:
		\begin{itemize}
			\item \texttt{Concepts}: characteristics of a statistical dataset (dimensions, attributes, measures)
			\item \texttt{Codelists}: description of a \textit{dimension} with a list of codes and labels
			\item \texttt{Datastructures}: description of the dataset structure
			\item \texttt{Data Structure Definitions (DSD)}: complete description of a data structure including the 3 previous types
		\end{itemize}
		\item \textit{Dataset} types:
		\begin{itemize}
			\item \texttt{GenericData}: generic SDMX data format
			\item \texttt{CompactData}: compacted data format
			\item \texttt{StructureSpecificData}: structure specific data format
			\item \texttt{UtilityData}: utility data format
			\item \texttt{MessageGroup}: specific message type developed to enable the exchange of several data or metadata messages of a single type in a unique SDMX-ML document. Currently enabled for datasets only.
		\end{itemize}
	\end{itemize}
  \end{frame}
 
  %rsdmx package - S4 model (Concepts)%
  %-----------------------------------%
  \begin{frame}[containsverbatim]
  	\frametitle{rsdmx}
  	\framesubtitle{Architecture - Object-oriented model (SDMX Concepts)}
  	\begin{itemize}
  		\item an \texttt{SDMXConcepts} object handles concepts either through \texttt{concepts} or \texttt{conceptSchemes} (depending on the SDMX version)
  		\item each concept is modeled with the \texttt{SDMXConcept} class
  	\end{itemize}
  	
	\centering
	\includegraphics[scale=0.2]{img/SDMXConcepts.png}
	
  \end{frame} 

  
  %rsdmx package - S4 model (Codelists)%
  %-----------------------------------%
  \begin{frame}[containsverbatim]
  	\frametitle{rsdmx}
  	\framesubtitle{Architecture - Object-oriented model (SDMX Codelists)}
  	\begin{itemize}
  		\item an \texttt{SDMXCodelists} object handles one or more codelists
  		\item each codelist is modeled with the \texttt{SDMXCodelist} class. It includes a list of \texttt{SDMXCode}
  	\end{itemize}
  	
  	\centering
	\includegraphics[scale=0.2]{img/SDMXCodelists.png}

  \end{frame} 
  
  %rsdmx package - S4 model (Data structures)%
  %------------------------------------------%
  \begin{frame}[containsverbatim]
  	\frametitle{rsdmx}
  	\framesubtitle{Architecture - Object-oriented model (SDMX Data structures / Key Families)}
  	\begin{itemize}
  		\item an \texttt{SDMXDataStructures} object handles one or more data structures (or key families)
  		\item each data structure is modeled with the \texttt{SDMXDataStructure} class. It includes a \texttt{SDMXComponents} object handling the dimensions, attributes, time dimension and measure
  	\end{itemize}
  	
  	\centering
	\includegraphics[scale=0.2]{img/SDMXDataStructures.png}

  \end{frame} 

  %rsdmx package - S4 model (DSD)%
  %------------------------------%
  \begin{frame}[containsverbatim]
  	\frametitle{rsdmx}
  	\framesubtitle{Architecture - Object-oriented model (SDMX Data structure Definition - DSD)}
  	an \texttt{SDMXDataStructures} object handles concepts, codelists and data structures.

  	\centering
	\includegraphics[scale=0.2]{img/SDMXDataStructureDefinition.png}

  \end{frame}


   %rsdmx package - object-oriented model%
  %-------------------------------------%
  \begin{frame}
  	\frametitle{rsdmx}
  	\framesubtitle{Architecture - readSDMX end-user function}
  	\texttt{readSDMX} is the main function of \texttt{rsdmx} package. The function will do the following:
  	\begin{itemize}
  		\item download the SDMX-ML document
  		\item determine the SDMX-ML message type and instantiate the corresponding \R \texttt{SDMX*} object
  		\item in case of \textit{Structure} message types, parse completely the document into a \texttt{S4} sub-model specific to the message type
  	\end{itemize}
  \end{frame}

  %rsdmx package - XML parsing technics %
  %-------------------------------------%
  \begin{frame}
  	\frametitle{rsdmx}
  	\framesubtitle{Architecture - XML Parsing technics \& strategies}
  	
  	2 different parsing technics:
  		\begin{itemize}
  			\item Initial and current technic: using \textbf{XPath}
  			\begin{itemize}
				\item requires loading the XML document tree in \R
				\item can cause R memory issues with large SDMX-ML documents
  			\end{itemize}
  			\vspace{0.1cm}
  			\item Alternative approach (in factory): using the \textbf{Simple API for XML (SAX)}
  			\begin{itemize}
  				\item does not require loading the XML document tree in \R
  				\item avoids R memory issues with large SDMX-ML documents
  			\end{itemize}
  			\vspace{0.1cm}
  			\item capacity to parse remote or local SDMX-ML files
  			\vspace{0.1cm}
  		\end{itemize}
  	\vspace{0.1cm}	
  	2 different parsing strategies:
  		\begin{itemize}
  			\item for \textit{Structure} types: when instantiating the \texttt{SDMX*} object (done by \texttt{readSDMX})
  			\item for \textit{Dataset} types: when coercing the \texttt{SDMX*} object to a \texttt{data.frame} (done by \texttt{as.data.frame})
  		\end{itemize}

  \end{frame}
 
\section{Using rsdmx} 
 
  %rsdmx package - Usage %
  %----------------------%
  %Installing rsdmx%
  \begin{frame}[containsverbatim]
  	\frametitle{rsdmx}
  	\framesubtitle{Usage - Installing rsdmx}
  	rsdmx can be installed:
  	\begin{itemize}
  		\item from CRAN
\begin{Schunk}
\begin{Sinput}
R> install.packages("rsdmx")
\end{Sinput}
\end{Schunk}
 	 
 	 	\vspace{0.1cm}
  		\item from Github (requires \texttt{devtools} package)
\begin{Schunk}
\begin{Sinput}
R> require(devtools)
R> install_github("eblondel/rsdmx")
\end{Sinput}
\end{Schunk}
  	\end{itemize}
	
	\vspace{0.1cm}
	Load rsdmx in R using:
\begin{Schunk}
\begin{Sinput}
     R> require(rsdmx)
\end{Sinput}
\end{Schunk}
		 
  \end{frame}
 
  %Datasets%
  \begin{frame}[containsverbatim]
  	\frametitle{rsdmx}
  	\framesubtitle{Usage - datasets}

Read a SDMX generic dataset in \R using the \textit{raw} approach.

\begin{Schunk}
\begin{Sinput}
R> url <- paste("http://data.fao.org/sdmx/repository/data/CROP_PRODUCTION/",
		".156.5312../FAO?startPeriod=2008&endPeriod=2008", sep="")
R> sdmxObj <- readSDMX(url)
R> class(sdmxObj)
\end{Sinput}
\begin{Soutput}
[1] "SDMXGenericData"
attr(,"package")
[1] "rsdmx"
\end{Soutput}
\end{Schunk}

Convert the \texttt{SDMXGenericData} into tabular data (\texttt{data.frame})

\begin{Schunk}
\begin{Sinput}
R> myData <- as.data.frame(sdmxObj)
R> head(myData)
\end{Sinput}
\begin{Soutput}
  FREQ REF_AREA INDICATOR COMMODITY DOMAIN UNITS UNIT_MULTIPLIER obsTime obsValue OBS_STATUS
1 YEAR      156      5312       515      Q    No            1000    2008     8832       <NA>
2 YEAR      156      5312       526      Q    No            1000    2008      450          E
3 YEAR      156      5312       367      Q    No            1000    2008      700          E
4 YEAR      156      5312       572      Q    No            1000    2008     4000          E
5 YEAR      156      5312        44      Q    No            1000    2008    67435       <NA>
6 YEAR      156      5312       414      Q    No            1000    2008      730          E
\end{Soutput}
\end{Schunk} 

  \end{frame}
  
   
  %Concepts%
  \begin{frame}[containsverbatim]
  	\frametitle{rsdmx}
  	\framesubtitle{Usage - concepts}
  	
Read a SDMX concepts document in \R

\begin{Schunk}
\begin{Sinput}
R> url <- paste("http://data.fao.org/sdmx/registry/conceptscheme/FAO/",
		"ALL/LATEST/?detail=full&references=none&version=2.1", sep="")
R> sdmxObj <- readSDMX(url)
R> class(sdmxObj)
\end{Sinput}
\begin{Soutput}
[1] "SDMXConcepts"
attr(,"package")
[1] "rsdmx"
\end{Soutput}
\end{Schunk}

Convert the \texttt{SDMXConcepts} into tabular data (\texttt{data.frame})

\begin{Schunk}
\begin{Sinput}
R> concepts <- as.data.frame(sdmxObj)
R> head(concepts[,c("id","en")])
\end{Sinput}
\begin{Soutput}
                id                         en
1        COMMODITY                  COMMODITY
2        INDICATOR                  INDICATOR
3         REF_AREA                   REF_AREA
4           DOMAIN                     DOMAIN
5     UNIT_MEASURE               UNIT_MEASURE
6             FREQ                       FREQ
7   FAO_MAJOR_AREA             FAO Major Area
8       UN_COUNTRY                 UN Country
9      ENVIRONMENT                Environment
10         SPECIES ASFIS Species Alpha 3 Code
11       OBS_VALUE                  OBS_VALUE
12      OBS_STATUS                 OBS_STATUS
13            UNIT                       Unit
14 UNIT_MULTIPLIER            Unit Multiplier
15            TIME                       Time

\end{Soutput}
\end{Schunk}   	
  	
  	
  \end{frame}
  
  %Codelists%
  \begin{frame}[containsverbatim]
  	\frametitle{rsdmx}
  	\framesubtitle{Usage - codelists}

Read a SDMX codelists document in \R

\begin{Schunk}
\begin{Sinput}
R> clUrl <- paste("http://data.fao.org/sdmx/registry/codelist/FAO/",
		"CL_FAO_MAJOR_AREA/0.1", sep="")
R> clObj <- readSDMX(clUrl)
R> class(clObj)
\end{Sinput}
\begin{Soutput}
[1] "SDMXCodelists"
attr(,"package")
[1] "rsdmx"
\end{Soutput}
\end{Schunk}

Convert the \texttt{SDMXCodelists} into tabular data (\texttt{data.frame})

\begin{Schunk}
\begin{Sinput}
R> codelist <- as.data.frame(clObj)
R> head(codelist[,c("id", "label.fr", "label.es")])
\end{Sinput}
\begin{Soutput}
  id                              label.fr                                label.es
1 01          Afrique - Eaux continentales            África - Aguas continentales
2 02 Amérique du Nord - Eaux continentales América del Norte - Aguas continentales
3 03  Amérique du Sud - Eaux continentales   América del Sur - Aguas continentales
4 04             Asie - Eaux continentales              Asia - Aguas continentales
5 05           Europe - Eaux continentales            Europa - Aguas continentales
6 06          Océanie - Eaux continentales           Oceanía - Aguas continentales
\end{Soutput}
\end{Schunk}   	
  	  	
 
  \end{frame}
  
 %DSD% 
 \begin{frame}[containsverbatim]
  	\frametitle{rsdmx}
  	\framesubtitle{Usage - Data structures}
  	
Read a SDMX Data Structure Definitions (DSD) document in \R

\begin{Schunk}
\begin{Sinput}
R> dsdUrl <- "http://stats.oecd.org/restsdmx/sdmx.ashx/GetDataStructure/TABLE1"
R> dsd <- readSDMX(dsdUrl)
R> class(dsd)
\end{Sinput}
\begin{Soutput}
[1] "SDMXDataStructureDefinition"
attr(,"package")
[1] "rsdmx"
\end{Soutput}
\end{Schunk}

Get the codelists contained in this DSD...

\begin{Schunk}
\begin{Sinput}
R> cls <- slot(dsd, "codelists")
R> codelists <- sapply(slot(cls, "codelists"), function(x) slot(x, "id"))
R> codelists
\end{Sinput}
\begin{Soutput}
 [1] "CL_TABLE1_OBS_STATUS"  "CL_TABLE1_DAC_DONOR"   "CL_TABLE1_DAC_PART"   
 [4] "CL_TABLE1_TRANSACTYPE" "CL_TABLE1_FLOWS"       "CL_TABLE1_DATATYPE"   
 [7] "CL_TABLE1_TIME"        "CL_TABLE1_UNIT"        "CL_TABLE1_POWERCODE"  "CL_TABLE1_TIME_FORMAT"
\end{Soutput}
\end{Schunk}

...and convert one codelist into tabular data (\texttt{data.frame})

\begin{Schunk}
\begin{Sinput}
R> cl <- as.data.frame(slot(dsd, "codelists"),codelistId = "CL_TABLE1_FLOWS")
R> cl
\end{Sinput}
\begin{Soutput}
    id          label.fr            label.en
1 1121   Versements Dons Grant Disbursements
2 1122  Versements PrOts  Loan Disbursements
3 1120  Versements Bruts Gross Disbursements
4 1130    Montants Recus    Amounts Received
5 1140   Versements Nets   Net Disbursements
6 1151  Engagements Dons   Grant Commitments
7 1152 Engagements PrOts    Loan Commitments
8 1150  Engagements nets         Commitments
\end{Soutput}
\end{Schunk}	
 
  \end{frame}
  
  
  %Datasets%
  \begin{frame}[containsverbatim]
  	\frametitle{rsdmx}
  	\framesubtitle{Usage - readSDMX with the helping approach}

\texttt{rsdmx} now brings the capacity to query data from a set of well-known data providers, still using the single \texttt{readSDMX} function. \texttt{rsdmx} embedds a list of SDMX service providers by default.

\vspace{\baselineskip}

The list of data providers "known" by \texttt{readSDMX} can be queried as follows:

\begin{Schunk}
\begin{Sinput}
R> providers <- getSDMXServiceProviders()
R> sapply(providers, slot, "agencyId")
\end{Sinput}
\begin{Soutput}
[1] "ECB"   "ESTAT" "OECD"  "FAO"   "ILO"  
\end{Soutput}
\end{Schunk}

The following example shows how to use \texttt{readSDMX} based on one of known data provider, OECD:

\begin{Schunk}
\begin{Sinput}
R> sdmx <- readSDMX(agencyId = "OECD", operation = "GetData", key = "MIG",
					filter = list("TOT", NULL, NULL),
					start = 2011, end = 2011)
\end{Sinput}
\end{Schunk}
  \end{frame}
  
  %Datasets%
  \begin{frame}[containsverbatim]
  	\frametitle{rsdmx}
  	\framesubtitle{Usage - readSDMX with the helping approach}

It is also possible to \textbf{add your own SDMX service provider}, and make it "known" by \texttt{readSDMX}!

\vspace{\baselineskip}

If you are interested, you can checkout the \texttt{rsdmx} documentation available \href{https://github.com/eblondel/rsdmx/wiki}{online}!

\vspace{\baselineskip}

If you want to register your SDMX service endpoint in the default list of providers, please contact me.

    \end{frame}

\section{Conclusions \& Perspectives}
  
  %rsdmx package - Success stories (Data sources) %
  %-----------------------------------------------%
  \begin{frame}
  	\frametitle{rsdmx}
  	\framesubtitle{Success stories - Variety of datasources}
  	Used on a variety of datasources:
  	\begin{itemize}
  		\item with \textbf{international} \& \textbf{regional} data sources: \href{http://data.un.org/Host.aspx?Content=API}{UN data portal}, \href{http://data.fao.org/sdmx/index.html}{UN Food \& Agriculture Organization (FAO)}, \href{http://www.ilo.org/ilostat/faces/home/statisticaldata/technical_page?_adf.ctrl-state=25zdozvi8_9&_afrLoop=1131342564621899}{UN International Labour Organization (ILO)} \href{http://stats.oecd.org/}{Organisation for Economic Co-operation and Development (OECD)}, \href{http://epp.eurostat.ec.europa.eu/portal/page/portal/sdmx_web_services/getting_started/rest_sdmx_2.1}{EUROSTAT}, \href{http://sdw.ecb.europa.eu/}{European Central Bank (ECB)}, \href{http://sdmxws.imf.org/Gateway2/}{International Monetary Fund (IMF)}, \href{http://databank.worldbank.org/data/home.aspx}{World Bank}
    	\item with \textbf{national} data sources: \href{http://govhack.abs.gov.au/}{Australian Bureau of Statistics (ABS)}, \href{https://www.ons.gov.uk/ons/apiservice/web/apiservice/home}{UK's Office of National Statistics (ONS)}, \href{http://www.bundesbank.de/}{Deutsche Bundesbank}, \href{http://www.bdm.insee.fr/bdm2/statique?page=sdmx}{INSEE (France)}, and more!!
  	\end{itemize} 	
  \end{frame}
    	
	%rsdmx package - Success stories (Projects) %
  %---------------------------------------------%
  \begin{frame}
  	\frametitle{rsdmx}
  	\framesubtitle{Success stories - Use in Projects}
  	Used in different web projects, such as:
  	\begin{itemize}
  		\item the \href{http://www.i-marine.eu/Pages/Home.aspx}{iMarine data e-infrastructure}: within R statistical analysis processings made available through Web Processing Services (WPS).
    	\item the \href{http://hackerspace.govhack.org/content/live-labour-force}{Live Labour Force project}: to allow reading SDMX datasets from the Australian Bureau of Statistics (ABS) portal (ABS.Stat).
    	\item the \href{http://syrtoproject.eu/}{SYRTO project}
  	\end{itemize}
  \end{frame}
  
  %rsdmx package - Perspectives %
  %-----------------------------%
  \begin{frame}
  	\frametitle{rsdmx}
  	\framesubtitle{Perspectives}
  	\begin{itemize}
  		\item enable the Simple API for XML (SAX) parsing technic for large datasets
  		\item improve the existing functionalities, \textit{e.g.} dataset time dimension format
  		\item support for additional SDMX document types and formats
  		\item extend the embedded list of SDMX service providers (\texttt{rsdmx} as web-service interface)
  		\item develop a generic SDMX web-client with the \texttt{shiny} R web-framework (\href{http://shiny.rstudio.com}{http://shiny.rstudio.com})
  	\end{itemize}
  \end{frame}
  
   %rsdmx package - Fundings %
  %--------------------------%
  \begin{frame}
  	\frametitle{rsdmx} 
   	\framesubtitle{Looking for Sponsors}	
  	\texttt{rsdmx} can play a \textbf{fundamental role} for \textbf{exploiting} and \textbf{co-analyzing statistics} from scattered data sources in \R.
  	
  	Until now, \texttt{rsdmx} was born from a \textbf{voluntary initiative}, and is now a \textbf{published library} with a \textbf{growing number of users}. To guarantee the sustainability of \texttt{rsdmx}, we are seeking for \href{https://github.com/eblondel/rsdmx/wiki\#package_development_funding}{\textbf{fundings}}, throught \textit{sponsoring} or \textit{donations} to:
  	\begin{itemize}
		\item \textbf{implement, test, validate and release improvements}
		\item \textbf{guarantee a quality maintenance} of the package
		\item \textbf{provide support} to users \& institutions that take advantage of \texttt{rsdmx}
  	\end{itemize}
  	
	If you are interesting in supporting \texttt{rsdmx}, please do not hesitate to contact \href{mailto:emmanuel.blondel1@gmail.com}{me}!  	
  	
  \end{frame}
  
  
  
  %rsdmx package - rsdmx on the web %
  %---------------------------------%
  \begin{frame}
  	\frametitle{rsdmx}
  	\framesubtitle{on the web}
  	\begin{itemize}
		\item on Github:
		\begin{itemize}
			\item source code: \href{https://github.com/eblondel/rsdmx}{https://github.com/eblondel/rsdmx}
			\vspace{0.1cm}
			\item online documentation: \href{https://github.com/eblondel/rsdmx/wiki}{https://github.com/eblondel/rsdmx/wiki}
			\vspace{0.1cm}
		\end{itemize}
		\item mailing list: \href{https://groups.google.com/forum/\#!forum/rsdmx}{rsdmx}			
		\vspace{0.1cm}
		\item on the Comprehensive R Archive Network (CRAN): \href{http://cran.r-project.org/package=rsdmx}{http://cran.r-project.org/package=rsdmx}
  	\end{itemize}  	
  \end{frame}  
  

\end{document}
